\documentclass[12pt, a4paper, french]{report}

% --- PACKAGES DE BASE ---
\usepackage[utf8]{inputenc}       % Encodage (standard)
\usepackage[T1]{fontenc}          % Polices de sortie
\usepackage[french]{babel}        % Langue française (guillemets, césure)
\usepackage{lmodern}              % Police plus lisible
\usepackage{microtype}            % Améliore l'apparence visuelle du texte

% --- PACKAGES MATHÉMATIQUES (Indispensables) ---
\usepackage{amsmath, amssymb, amsthm, amsfonts} % La base du LaTeX mathématique
\usepackage{mathtools}             % Extensions pour amsmath
\usepackage{mathrsfs}              % Lettres calligraphiques (\mathscr{L})
\usepackage{bm}                    % Symboles en gras (\bm{x})

% --- CONFIGURATION DES THÉORÈMES ---
\newtheorem{theorem}{Théorème}[chapter]
\newtheorem{lemma}[theorem]{Lemme}
\newtheorem{proposition}[theorem]{Proposition}
\newtheorem{corollary}[theorem]{Corollaire}

\theoremstyle{definition}
\newtheorem{definition}[theorem]{Définition}
\newtheorem{example}[theorem]{Exemple}
\newtheorem{remark}[theorem]{Remarque}

% --- MISE EN PAGE ---
\usepackage[margin=2.5cm]{geometry} % Marges
\usepackage{hyperref}               % Liens cliquables (sommaire, citations)
\hypersetup{colorlinks=true, linkcolor=blue, citecolor=red}

% --- TES MACROS (Raccourcis) ---
\newcommand{\R}{\mathbb{R}}
\newcommand{\N}{\mathbb{N}}
\newcommand{\C}{\mathbb{C}}
\newcommand{\diff}{\mathrm{d}}     % Pour un joli 'd' de différentielle

% --- INFOS DU MÉMOIRE ---
\title{Titre de ton Mémoire de Recherche}
\author{Gabriel de Boisgrollier}
\date{\today}

\begin{document}

\maketitle

\tableofcontents

% --- INCLUSION DES CHAPITRES ---
% Git va adorer cette séparation par fichiers !
\chapter{Introduction}
\cite{BOT_2024}
Let us define the parabolic differential operator of second order :
$$L \coloneq \partial_0 - \sum_{i=1}^{d}\partial_i^2.$$
Here, $x_0$ is the time variable and $\{x_i\}_{i=1,\dots,d}$ are the space variable.
We are insterested in the equation 
\begin{equation}
    L\phi - \lambda \phi^3 = \xi \quad \text{on }\R^d
    \label{eq:phi4}
\end{equation}
where $\lambda$ is a real parameter, and $\xi$ a space time function/distribution.
This equation presents a non-linearity (when $\lambda \neq 0$). We are insterested in the cases where $\xi$ is rough, and thus the equation does not have a rigourous meaning a priori.
Indeed, if $\xi$ is a distribution so rough that the solution is also not better that a distribution, then $\phi^3$ is ill-defined.
For example, if $\lambda = 0$, $d = 0$ and $\xi = \frac{d^2B}{dx_0^2}$, then the solution $\Pi_0$ is given by $\frac{dB}{dx_0}$ which is a genuine distribution of regularity just worse than $- \frac{1}{2}$.



Hello world 

% --- BIBLIOGRAPHIE ---
% \bibliographystyle{alpha}
% \bibliography{biblio.bib}

\end{document}